% !TeX root = ../complete.tex
\chapter{تأخیر و اندازه‌گذاری گیت‌ها با \lr{Logical Effort}}
\label{chap:delay}

وقتی ورودی یک گیت تغییر می‌کند، خروجی چه زمانی تغییر می‌کند؟ برای
پاسخ، ابتدا زمان را روی شکل موج اندازه می‌گیریم. سپس از روی مسیر
شارژ یا تخلیه و ظرفیت گره‌ها، تأخیر را تخمین می‌زنیم. در ادامه فصل،
همین مدل پایه را برای مقایسه گیت‌ها، انتخاب تعداد مراحل و اندازه‌گذاری
گیت‌های یک مسیر به کار می‌بریم.

پیش از شروع، \NotesTopic{sec:book-cmos}{مسیرهای هدایت در گیت‌های
\lr{CMOS}}، \NotesTopic{eq:foundation-rc}{پاسخ مدار \lr{RC} تک‌خازنی}
و \NotesTopic{sec:book-cv}{مبدأ ظرفیت‌های ترانزیستور} را مرور کنید.
در خواندن نخست، از روی هر مدار مشخص کنید کدام گره تغییر ولتاژ می‌دهد،
بار آن از کجا می‌آید و کدام ترانزیستور جریان شارژ یا تخلیه را فراهم می‌کند.
هنگام مرور، از شکل موج به مدل مدار بروید و سپس بررسی کنید که مقدار
به‌دست‌آمده با افزایش بار یا کاهش قدرت راه‌انداز چگونه تغییر می‌کند.

% !TeX root = ../complete.tex
\section{شکل موج، معیارهای زمانی و پاسخ گذرا}
\label{sec:delay-transient}

در \NotesTopic{sec:dc-transfer}{مشخصه انتقال \lr{DC}}، خروجی ماندگار
را برای هر ولتاژ ورودی تعیین کردیم. در تحلیل گذرا، زمان رسیدن به سطح
منطقی معتبر نیز اهمیت دارد. گیتی که گره خروجی را شارژ یا تخلیه می‌کند،
\lr{Driver} است. گیت‌های متصل به آن گره و سیم رابط، بخشی از \lr{Load}
آن هستند. ظرفیت‌های خود \lr{Driver} نیز در تعیین جریان لازم برای تغییر
ولتاژ خروجی سهم دارند \cite[بخش‌های \secnum{۴.۱.۱} و
\secnum{۴.۲}، صفحات ۱۴۱ تا ۱۴۴]{weste}.

\subsection{اندازه‌گیری زمان روی شکل موج}
\label{sec:delay-waveform}

برای شکل موجی که میان زمین و \(V_{DD}\) تغییر می‌کند، مرجع معمول
اندازه‌گیری تأخیر، \(V_{DD}/2\) در هر دو شکل موج است. فاصله زمانی
بین عبور ورودی و عبور خروجی از این مرجع، تأخیر همان گذار است.
برای شرایط مشخص، بیشترینِ تأخیرهای معتبر را \lr{Propagation Delay}
با نماد \(t_{\mathrm{pd}}\) و کمترین را \lr{Contamination Delay}
با نماد \(t_{\mathrm{cd}}\) می‌نامیم. این دو حد، دیرترین و زودترین
عبور خروجی از مرجع را مشخص می‌کنند. در برآورد دقیق‌تر می‌توان حدها
را برای هر ورودی و برای گذار صعودی یا نزولی خروجی جداگانه مشخص کرد.
برای نمونه، \(t_{\mathrm{pdf}}\) حد بالای تأخیر گذار نزولی خروجی است
\cite[بخش \secnum{۴.۱.۱} و شکل \secnum{۴.۱}، صفحات ۱۴۱ و ۱۴۲]{weste}.

\lr{Rise Time} با نماد \(t_{\mathrm{r}}\) زمان عبور یک شکل موج از
۲۰ تا ۸۰ درصد دامنه تغییر و \lr{Fall Time} با نماد \(t_{\mathrm{f}}\)
زمان عبور آن از ۸۰ تا ۲۰ درصد است. این بازه‌ها قرارداد کتاب‌اند.
میانگین \((t_{\mathrm{r}}+t_{\mathrm{f}})/2\) را کتاب به‌عنوان
\lr{Edge Rate} با واحد زمان معرفی می‌کند. در شکل
\ref{fig:delay-waveform}، پیکان بالا فاصله میان دو شکل موج و پیکان
پایین فاصله میان دو نقطه از همان شکل موج خروجی را نشان می‌دهد
\cite[بخش \secnum{۴.۱.۱}، صفحه ۱۴۱]{weste}.

\begin{figure}[htbp]
  \centering
  \begin{latin}
  \begin{tikzpicture}[x=1.52cm,y=2.35cm]
    \draw[->] (0,0) -- (7,0) node[right] {\(t\)};
    \draw[->] (0,0) -- (0,1.25) node[above] {\(V/V_{DD}\)};
    \foreach \level/\ticklabel in {0.2/۰٫۲,0.5/۰٫۵,0.8/۰٫۸,1/۱} {
      \draw[gray!60,thin,dotted] (0,\level) -- (6.65,\level);
      \node[left] at (0,\level) {\fontspec[Path=C:/Windows/Fonts/]{arial.ttf}\ticklabel};
    }
    \draw[very thick] (0,0) -- (0.5,0) -- (2.5,1) -- (6.65,1);
    \draw[very thick,blue!70!black] (0,1) -- (3,1) -- (6,0) -- (6.65,0);
    \node[above] at (5.7,1) {\(V_A\)};
    \node[above,blue!70!black] at (0.85,1) {\(V_Y\)};
    \draw[densely dashed] (1.5,0.5) -- (1.5,1.22);
    \draw[densely dashed] (4.5,0.5) -- (4.5,1.22);
    \fill (1.5,0.5) circle (1.6pt) (4.5,0.5) circle (1.6pt);
    \draw[<->] (1.5,1.18) -- (4.5,1.18) node[midway,above] {\(\Delta t_{50}\)};
    \draw[densely dashed] (3.6,0.8) -- (3.6,-0.21);
    \draw[densely dashed] (5.4,0.2) -- (5.4,-0.21);
    \fill[blue!70!black] (3.6,0.8) circle (1.6pt) (5.4,0.2) circle (1.6pt);
    \draw[<->] (3.6,-0.17) -- (5.4,-0.17) node[midway,below] {\(t_{\mathrm{f}}\)};
  \end{tikzpicture}
  \end{latin}
  \caption{شکل موج شماتیک ورودی \(V_A\) و خروجی \(V_Y\) یک \lr{Inverter}.
  \(\Delta t_{50}\) تأخیر یک گذار و \(t_{\mathrm{f}}\) زمان افت همان خروجی است.
  قرارداد اندازه‌گیری از شکل \secnum{۴.۱} کتاب \cite{weste} پیروی می‌کند.}
  \label{fig:delay-waveform}
\end{figure}

برای مثال، اگر ورودی در \qty{20}{\pico\second} و خروجی نزولی در
\qty{55}{\pico\second} از مرجع نصف تغذیه بگذرند، تأخیر این گذار
\qty{35}{\pico\second} است. اگر خروجی در
\qty{43}{\pico\second} از ۸۰ درصد و در \qty{70}{\pico\second}
از ۲۰ درصد عبور کند، زمان افت آن \qty{27}{\pico\second} است.
برای یافتن \(t_{\mathrm{pd}}\) و \(t_{\mathrm{cd}}\) گیت، دامنه
ورودی‌ها، گذارها، بار و شرایط کاری مورد بررسی باید مشخص باشند.

\subsection{تأخیر گیت و زمان رسیدن در شبکه منطقی}
\label{sec:delay-arrival-time}

زمانی را که گذار به گره \(i\) می‌رسد با \(a_i\) نمایش می‌دهیم.
در مدل مقدماتی یک شبکه ترکیبی، \(a_i\) را با دیرترین زمان رسیدن
به ورودی‌های گیتِ راه‌انداز گره \(i\) و حد بالای تأخیر همان گیت
تخمین می‌زنیم. مجموعه این ورودی‌ها را \(\operatorname{fanin}(i)\)
و حد بالای مشترک تأخیر گیت را \(t_{\mathrm{pd},i}\) می‌نامیم
\cite[رابطه \secnum{۴.۱}، صفحه ۱۴۲]{weste}.
\begin{equation}
  a_i=\max_{j\in\operatorname{fanin}(i)} a_j+t_{\mathrm{pd},i}.
  \label{eq:delay-arrival}
\end{equation}
اگر زمان رسیدن دو ورودی \qty{25}{\pico\second} و
\qty{40}{\pico\second} و تأخیر گیت \qty{30}{\pico\second} باشد،
مدل، زمان رسیدن خروجی را \qty{70}{\pico\second} می‌دهد. اگر
خروجی تا \qty{85}{\pico\second} لازم باشد، \lr{Slack}، یعنی
زمان لازم منهای زمان رسیدن، \qty{15}{\pico\second} است. در تحلیل
دقیق‌تر، تأخیر هر ورودی، جهت گذار و امکان ایجاد گذار در خروجی
بررسی می‌شود \cite[بخش \secnum{۴.۱.۱}، صفحه ۱۴۲]{weste}.

مسیرهایی که مهلت زمانی را محدود می‌کنند \lr{Critical Path} هستند.
می‌توان تأخیر را در سطح \lr{Microarchitecture} با تغییر تقسیم کار میان
مراحل، در سطح منطق با انتخاب توپولوژی، در سطح مدار با اندازه‌گذاری
ترانزیستورها و در سطح \lr{Layout} با کاهش بار سیم و اتصال‌ها تغییر داد.
این فصل بر تعداد مراحل، نوع گیت و اندازه‌گذاری در سطح منطق و مدار
تمرکز دارد \cite[بخش \secnum{۴.۱.۲}، صفحات ۱۴۲ و ۱۴۳]{weste}.

\subsection{چرا تغییر ولتاژ خروجی زمان می‌برد؟}
\label{sec:delay-inverter-transient}

خروجی \(Y\) یک \lr{Inverter} را در نظر بگیرید که با یک سیم کوتاه،
ورودی گیت بعدی را راه‌اندازی می‌کند. در تقریب ظرفیت متمرکز، مقاومت
سیم ناچیز فرض می‌شود و ظرفیت سیم، ظرفیت‌های \lr{Diffusion} متصل به
\(Y\) و ظرفیت ورودی گیت بعدی با هم جمع می‌شوند. ظرفیت معادل حاصل را
\(C_{\mathrm{out}}\) می‌نامیم. ظرفیت یک پایانه متصل به گره دیگری
به‌صورت مستقیم در این جمع وارد نمی‌شود. برای مثال، ظرفیت
\lr{Source-to-Body} با دو پایانه در ولتاژ ثابت، در این گذار بار
اضافی از خروجی نمی‌گیرد \cite[بخش \secnum{۴.۲} و شکل
\secnum{۴.۳}، صفحات ۱۴۳ و ۱۴۴]{weste}.

ورودی را ابتدا صفر بگیرید. خروجی در \(V_{DD}\) است. پس از گذار
پله‌ای ورودی به \(V_{DD}\)، در مدل ایدئال \lr{pMOS} خاموش و
\lr{nMOS} روشن است. اگر \(I_n\) بزرگی جریان خروجی به زمین و
\(C_{\mathrm{out}}\) ثابت باشد، تعادل جریان گره می‌دهد
\begin{equation}
  C_{\mathrm{out}}\frac{\mathrm{d}V_Y}{\mathrm{d}t}=-I_n.
  \label{eq:delay-step-kcl}
\end{equation}
علامت منفی یعنی جریان خروجی، بار خازن و به‌دنبال آن \(V_Y\) را کاهش
می‌دهد. واحد هر دو طرف آمپر است. این رابطه صورت لحظه‌ای همان
\NotesTopic{eq:impl-charge-delay}{رابطه بار، جریان و زمان} است.

در مدل \NotesTopic{sec:book-mos-current}{\lr{Long-channel}}، اگر
\(V_{DD}>V_{tn}\)، اضافه‌ولتاژ \(V_{GT}=V_{DD}-V_{tn}\) است.
وقتی \(V_Y\) از \(V_{GT}\) بیشتر است، ترانزیستور در \lr{Saturation}
قرار دارد. پس از رسیدن خروجی به \(V_{GT}\)، ترانزیستور وارد
\lr{Linear Region} می‌شود. با صرف‌نظر از اثرهای غیرایده‌آل،
\(\beta_n=\mu_n C_{\mathrm{ox}}W/L\) ضریب فناوری و هندسه است و
رابطه جریان را می‌توان با شرط هر ناحیه نوشت \cite[رابطه
\secnum{۴.۴}، صفحه ۱۴۴]{weste}.
\begin{equation}
  \label{eq:delay-step-current}
  I_n=
  \begin{cases}
    \dfrac{\beta_n}{2}V_{GT}^{2},
      & V_Y\ge V_{GT},\\[3pt]
    \beta_n\left(V_{GT}V_Y-\dfrac{V_Y^2}{2}\right),
      & 0\le V_Y<V_{GT}.
  \end{cases}
\end{equation}
دو عبارت در مرز \(V_Y=V_{GT}\) برابرند. در ناحیه اول، جریان
ثابت است و از رابطه \eqref{eq:delay-step-kcl}، ولتاژ با شیب ثابت
\(-\beta_n V_{GT}^2/(2C_{\mathrm{out}})\) پایین می‌آید. با ورود به
ناحیه دوم، جریان با کاهش \(V_Y\) کم می‌شود. اگر
\(V_{DD}/2<V_{GT}\) باشد، رسیدن خروجی به مرجع نصف تغذیه پس از
ورود به ناحیه دوم رخ می‌دهد و باید سهم زمانی هر دو ناحیه را حساب کرد.
مدل \lr{RC} در بخش \ref{sec:delay-rc-model} راهی برای تخمین زمان
عبور بدون حل جداگانه معادله هر ناحیه فراهم می‌کند.

وقتی ورودی به‌صورت \lr{Ramp} بالا می‌رود، جریان هر دو ترانزیستور
به \(V_A(t)\) وابسته است. برای \(V_{tn}+|V_{tp}|<V_{DD}\)، ابتدا
\lr{nMOS} خاموش است. میان \(V_{tn}\) و \(V_{DD}-|V_{tp}|\)،
هر دو ترانزیستور هدایت می‌کنند و جریان خالص خازن \(I_p-I_n\) است.
در این عبارت، \(I_p\) بزرگی جریان از تغذیه به خروجی است. بالاتر از
این بازه، \lr{pMOS} خاموش می‌شود و تخلیه از مسیر \lr{nMOS} ادامه
می‌یابد. در نتیجه، هنگام مقایسه تأخیرها باید شکل موج ورودی و بار
خروجی را مشخص کرد \cite[بخش \secnum{۴.۲} و جدول
\secnum{۴.۱}، صفحات ۱۴۵ و ۱۴۶]{weste}.

% !TeX root = ../complete.tex
\section{مدل تأخیر \lr{RC} و \lr{Elmore Delay}}
\label{sec:delay-rc-model}

برای تخمین دستی تأخیر، ترانزیستور روشن را با یک مقاومت و ظرفیت‌های
وابسته به ولتاژ را با مقادیر نماینده در بازه گذار مدل می‌کنیم.
هدف این مدل، نشان‌دادن اثر مسیر هدایت و بار روی زمان کلیدزنی است.
دقت تخمین به گذار، مقدار مقاومت و ظرفیت‌های انتخاب‌شده وابسته است
\cite[بخش \secnum{۴.۳}، صفحات ۱۴۶ تا ۱۵۵]{weste}.

\subsection{مقاومت، ظرفیت و انتخاب مسیر هدایت}
\label{sec:delay-rc-sizing}

در مدل کتاب، ترانزیستور خاموش مسیر جریان را قطع می‌کند. ترانزیستور
روشن مانند کلیدی بسته با مقاومت مؤثر در مسیر هدایت عمل می‌کند.
اگر عرض افزاره \(k\) برابر عرض واحد باشد و طول کانال و شرایط مدل
ثابت بمانند، تقریب می‌زنیم که جریان \(k\) برابر و مقاومت \(1/k\)
برابر شود. مقاومت واحد \lr{nMOS} را \(R\) می‌نامیم. در قرارداد
آموزشی کتاب، مقاومت واحد \lr{pMOS} برابر \(2R\) است. بنابراین
مقاومت مدل برای \lr{nMOS} و \lr{pMOS} با ضریب عرض \(k\)
به‌ترتیب \(R/k\) و \(2R/k\) است. ضریب ۲ فرض مثال است و به فناوری
وابسته است \cite[بخش \secnum{۴.۳.۱}، صفحه ۱۴۶]{weste}.

اگر \(C\) ظرفیت ورودی ترانزیستور واحد باشد، ظرفیت ورودی ترانزیستور
با ضریب عرض \(k\) را \(kC\) می‌گیریم. هر ناحیه تماس‌دار
\lr{Diffusion} در افزاره واحد نیز در تقریب آموزشی ظرفیتی حدود
\(C\) دارد. ظرفیت‌های روی یک گره جمع می‌شوند. جمله‌های ظرفیت، پس
از انتخاب مسیر شارژ یا تخلیه، باید از روی اتصال‌های همان مدار شمرده
شوند \cite[بخش‌های \secnum{۴.۳.۲} و \secnum{۴.۳.۳}، صفحه ۱۴۷]{weste}.

برای \lr{Inverter} مرجع با عرض‌های نسبی ۱ برای \lr{nMOS} و ۲
برای \lr{pMOS}، ورودی مشترک ظرفیت \(C+2C=3C\) دارد. در افت
خروجی، مسیر تخلیه از \lr{nMOS} می‌گذرد و در صعود خروجی، مسیر
شارژ از \lr{pMOS} می‌گذرد. مقاومت هر دو مسیر در این مدل \(R\)
است، زیرا مقاومت \lr{pMOS} با عرض ۲ برابر \(2R/2\) می‌شود.
دو ناحیه \lr{Drain} متصل به خروجی در تقریب ساده ظرفیت \(3C\)
دارند. اگر خروجی، ورودی یک \lr{Inverter} هم‌اندازه را نیز
راه‌اندازی کند و ظرفیت سیم قابل صرف‌نظر باشد، بار گره \(6C\)
است \cite[شکل \secnum{۴.۶}، صفحه ۱۴۷]{weste}. در گیت‌هایی
با ترانزیستورهای سری، مقاومت مسیر از جمع مقاومت ترانزیستورهای
روشن به دست می‌آید. گره‌های میان آن‌ها نیز می‌توانند ظرفیت مؤثر
در گذار داشته باشند.

\subsection{پاسخ تک‌خازنی و قرارداد زمان نصف تغذیه}
\label{sec:delay-rc-half-time}

ابتدا یک مقاومت خطی \(R_{\mathrm{phys}}\) و یک خازن \(C_\Sigma\)
را در نظر بگیرید. با اعمال پله ایدئال در \(t=0\)،
\NotesTopic{eq:foundation-rc}{پاسخ \lr{RC} فصل پیش‌نیاز}
برای شارژ از صفر به \(V_{DD}\) قابل استفاده است. اگر خازن ابتدا
در \(V_{DD}\) باشد و از راه مقاومت به زمین تخلیه شود، در معادله
گره به‌جای \(V_{DD}-v(t)\)، اختلاف ولتاژ تخلیه \(v(t)\) قرار
می‌گیرد و جریان از خازن خارج می‌شود. با تعریف ثابت زمانی
\(\tau_{\mathrm{RC}}=R_{\mathrm{phys}}C_\Sigma\)، دو پاسخ چنین‌اند
\cite[بخش \secnum{۴.۳.۴}، صفحات ۱۴۸ و ۱۴۹]{weste}.
\begin{equation}
  \begin{aligned}
    v_{\mathrm{rise}}(t)&=V_{DD}\bigl(1-e^{-t/\tau_{\mathrm{RC}}}\bigr),\\
    v_{\mathrm{fall}}(t)&=V_{DD}e^{-t/\tau_{\mathrm{RC}}},
      \qquad t\ge 0.
  \end{aligned}
  \label{eq:delay-rc-step}
\end{equation}
در \(t=0\)، ولتاژ شارژ صفر و ولتاژ تخلیه \(V_{DD}\) است.
در هر دو پاسخ، شرط \(v(t_{50})=V_{DD}/2\) نتیجه می‌دهد که
\(e^{-t_{50}/\tau_{\mathrm{RC}}}=1/2\). با گرفتن لگاریتم طبیعی
از دو طرف و استفاده از \(\ln(1/2)=-\ln 2\) داریم
\begin{equation}
  t_{50}=R_{\mathrm{phys}}C_\Sigma\ln 2.
  \label{eq:delay-rc-physical-half-time}
\end{equation}
برای نمونه، \(R_{\mathrm{phys}}=\qty{2}{\kilo\ohm}\) و
\(C_\Sigma=\qty{10}{\femto\farad}\) ثابت زمانی
\qty{20}{\pico\second} و \(t_{50}\approx\qty{13.9}{\pico\second}\)
می‌دهند. این اعداد فرضی‌اند. اگر بار دو برابر شود و مقاومت تغییر
نکند، زمان رسیدن به نصف تغذیه نیز دو برابر می‌شود.

کتاب ضریب \(\ln 2\) را در مقاومت مؤثرِ مدل زمان جذب می‌کند.
اگر \(R_{\mathrm{eff}}=R_{\mathrm{phys}}\ln 2\) انتخاب شود،
رابطه \eqref{eq:delay-rc-physical-half-time} به
\(t_{50}=R_{\mathrm{eff}}C_\Sigma\) تبدیل می‌شود. کتاب سپس
برای ساده‌نویسی \(R_{\mathrm{eff}}\) را با \(R\) می‌نویسد.
در این قرارداد، برای \lr{Inverter} مرجعِ بخش
\ref{sec:delay-rc-sizing} با بار کل \(6C\)، تخمین تأخیر
\(6RC\) است و ضریب \(\ln 2\) پیش‌تر در \(R\) منظور شده
است \cite[رابطه \secnum{۴.۹} و شکل \secnum{۴.۶}، صفحات
۱۴۷ و ۱۴۹]{weste}.

\subsection{از یک گره تا درخت \lr{RC}}
\label{sec:delay-elmore}

وقتی سیم یا شاخه‌های گیت مقاومت قابل‌توجه دارند، گره‌ها می‌توانند با
سرعت‌های متفاوت شارژ شوند. شکل \ref{fig:delay-elmore-tree} نمونه‌ای
از درخت \lr{RC} با یک مسیر مشترک و دو شاخه است. درخت حلقه ندارد،
پس از منبع تا هر گره تنها یک مسیر وجود دارد. در شکل، \(Y\) خروجی
موردنظر و \(Z\) انتهای شاخه دیگر است. \(C_X\)، \(C_Y\) و
\(C_Z\) به‌ترتیب ظرفیت‌های گره‌های \(X\)، \(Y\) و \(Z\) هستند.

\begin{figure}[htbp]
  \centering
  \begin{latin}
  \begin{circuitikz}
    \draw (0,0) node[ocirc,label=above:\(V_s(t)\)] {}
      to[R,l=\(R_0\)] (3,0) node[circ,label=above left:\(X\)] {}
      to[R,l=\(R_1\)] (6,0) node[circ,label=right:\(Y\)] {}
      to[C,l=\(C_Y\)] (6,-1.8) node[ground] {};
    \draw (3,0) to[C,l=\(C_X\)] (3,-1.8) node[ground] {};
    \draw (3,0) -- (3,1.45)
      to[R,l=\(R_2\)] (8.5,1.45) node[circ,label=right:\(Z\)] {}
      to[C,l=\(C_Z\)] (8.5,-1.8) node[ground] {};
  \end{circuitikz}
  \end{latin}
  \caption{درخت \lr{RC} برای محاسبه تأخیر از منبع \(V_s(t)\) تا \(Y\).
  بار شاخه \(Z\) نیز از مقاومت مشترک \(R_0\) شارژ می‌شود.
  شکل برای مثال آموزشی و بر اساس روش بخش \secnum{۴.۳.۵} کتاب
  رسم شده است \cite{weste}.}
  \label{fig:delay-elmore-tree}
\end{figure}

برای گره هدف \(q\)، از بین مقاومت‌های مسیر منبع تا گره \(i\)،
فقط مقاومت‌هایی را جمع می‌کنیم که در مسیر منبع تا \(q\) نیز
هستند. این جمع \(R_{\mathrm{common}}(i,q)\) است. اگر \(C_i\)
ظرفیت گره \(i\) باشد، ثابت زمانی \lr{Elmore} برای گره هدف از
جمع وزن‌دار زیر به دست می‌آید \cite[روابط \secnum{۴.۱۳}
و \secnum{۴.۱۴}، صفحه ۱۵۰]{weste}.
\begin{equation}
  \tau_{\mathrm{E},q}=\sum_i R_{\mathrm{common}}(i,q)C_i.
  \label{eq:delay-elmore-common-path}
\end{equation}
برای فهم وزن‌ها، فرض کنید منبع از صفر به یک ولتاژ ثابت پله می‌زند
و همه خازن‌ها ابتدا تخلیه‌اند. بار نهایی هر خازن برابر ظرفیتِ آن
ضرب‌در تغییر ولتاژ منبع است. تمام این بار از مقاومت‌های میان
منبع و آن خازن عبور می‌کند. از سوی دیگر، اختلاف ولتاژ لحظه‌ای
بین منبع و گره هدف فقط از افت مقاومت‌های مسیر همان گره تشکیل
می‌شود. بنابراین برای محاسبه مساحت زیر منحنی اختلاف ولتاژ منبع
و گره هدف، سهم هر بار تنها از مقاومت‌های مشترک وارد می‌شود.
تقسیم این مساحت بر دامنه پله منبع، رابطه
\eqref{eq:delay-elmore-common-path} را با واحد زمان می‌دهد.

برای خروجی \(Y\) در شکل \ref{fig:delay-elmore-tree}، مسیر منبع
تا \(Y\) شامل \(R_0\) و \(R_1\) است. بارهای \(X\) و \(Z\)
از \(R_0\) سهم می‌گیرند و بار \(Y\) از هر دو مقاومت. در نتیجه
\begin{equation}
  \tau_{\mathrm{E},Y}
    =R_0C_X+(R_0+R_1)C_Y+R_0C_Z
    =R_0(C_X+C_Y+C_Z)+R_1C_Y.
  \label{eq:delay-elmore-branch-example}
\end{equation}
مقاومت \(R_2\) در مسیر منبع تا \(Y\) قرار ندارد. مقدار آن
می‌تواند شکل موج واقعی را تغییر دهد، ولی در این جمع ظاهر
نمی‌شود. با حذف شاخه \(Z\)، رابطه به فرم مدار دوگرهی کتاب
کاهش می‌یابد \cite[شکل \secnum{۴.۱۰} و مثال
\secnum{۴.۳}، صفحات ۱۴۹ و ۱۵۰]{weste}.

در یک مثال فرضی، \(R_0=\qty{1}{\kilo\ohm}\)،
\(R_1=\qty{2}{\kilo\ohm}\)، \(R_2=\qty{3}{\kilo\ohm}\) و
\(C_X\)، \(C_Y\) و \(C_Z\) به‌ترتیب \qty{10}{\femto\farad}،
\qty{20}{\femto\farad} و \qty{30}{\femto\farad} هستند. با
توجه به اینکه حاصل‌ضرب یک \unit{\kilo\ohm} در یک
\unit{\femto\farad} برابر یک \unit{\pico\second} است، رابطه
\eqref{eq:delay-elmore-branch-example} می‌دهد
\[
  \tau_{\mathrm{E},Y}
    =[1(10+20+30)+2(20)]\,\unit{\pico\second}
    =\qty{100}{\pico\second}.
\]
اگر پاسخ را با یک منحنی نمایی با همین ثابت زمانی تقریب بزنیم،
\(t_{50}\approx\tau_{\mathrm{E},Y}\ln 2\) و در این مثال
\(t_{50}\approx\qty{69.3}{\pico\second}\) است. این عدد
از تقریب تک‌نمایی به دست آمده و شکل موج درخت به‌طور کامل حل
نشده است. کتاب در محاسبه‌های دستی، ضریب \(\ln 2\) را در
مقاومت‌ها منظور می‌کند و همان جمع وزن‌دار را برای تخمین تأخیر
به کار می‌برد \cite[بخش‌های \secnum{۴.۳.۴} و
\secnum{۴.۳.۵}، صفحات ۱۴۹ و ۱۵۰]{weste}.

\subsection{اثر \lr{Layout} و تعیین مقاومت مدل}
\label{sec:delay-rc-calibration}

فرض تناسب ظرفیت \lr{Diffusion} با عرض، برای تخمین اولیه مفید
است. پس از تعیین \lr{Layout}، باید ناحیه‌های مشترک و اتصال‌های
واقعی را شمرد. در نمونه \lr{NAND} سه‌ورودی کتاب، دو \lr{Drain}
در شبکه \lr{pMOS} یک ناحیه تماس‌دار مشترک دارند. ظرفیت تخمینی
خروجی در آن طرح از \(9C\) به \(7C\) می‌رسد. تقسیم یک ترانزیستور
پهن به چند شاخه موازی نیز، بسته به چیدمان و اشتراک ناحیه‌ها،
می‌تواند ظرفیت را تغییر دهد \cite[بخش \secnum{۴.۳.۶} و
شکل‌های \secnum{۴.۱۷} و \secnum{۴.۱۸}، صفحات ۱۵۳ و ۱۵۴]{weste}.

برای تخمین مقاومت مؤثر از جریان، تخلیه خروجی از \(V_{DD}\) تا
\(V_{DD}/2\) را در نظر بگیرید. اگر \(I_{\mathrm{eff}}\) جریان
نماینده در همین بازه باشد، باید بار
\(C_\Sigma V_{DD}/2\) جابه‌جا شود. تقسیم بار بر جریان، زمان
تخمینی را می‌دهد. یکسان‌گرفتن این زمان با
\(R_{\mathrm{eff}}C_\Sigma\) نتیجه می‌دهد
\begin{equation}
  \label{eq:delay-rc-current-calibration}
  t_{50}\approx\frac{C_\Sigma V_{DD}}{2I_{\mathrm{eff}}},
  \qquad R_{\mathrm{eff}}\approx\frac{V_{DD}}{2I_{\mathrm{eff}}}.
\end{equation}
در تخمین \lr{Ramp} کتاب، جریان نماینده را از میانگین جریان
ابتدا و انتهای بازه می‌گیرند. شرایط بایاس، شیب ورودی و جهت
گذار در انتخاب مقدار مؤثر نقش دارند. مقدار دقیق‌تر را می‌توان
با مدل افزاره و شبیه‌سازیِ همان گذار به دست آورد
\cite[بخش \secnum{۴.۳.۷} و روابط \secnum{۴.۱۷} تا
\secnum{۴.۱۹}، صفحات ۱۵۴ و ۱۵۵]{weste}.

مدل \lr{RC} نشان می‌دهد چرا بار بیشتر، مسیر مقاوم‌تر و ظرفیت
شاخه‌های مشترک می‌توانند تأخیر را افزایش دهند. در مدل خطی تأخیر،
همین رفتار را برای مقایسه گیت‌ها با دو سهمِ وابسته به بار و
وابسته به ظرفیت‌های درونی بیان می‌کنیم.

% !TeX root = ../complete.tex
\section{مدل خطی تأخیر و \lr{Logical Effort} تک‌گیت}
\label{sec:delay-linear-model}

در \hyperref[sec:delay-rc-model]{مدل \lr{RC}}، افزایش ظرفیت بارِ
یک گیت با مقاومت ثابت، تأخیر آن را افزایش می‌دهد. مدل خطی همین
رفتار را با شیب و عرض از مبدأ یک خط بیان می‌کند. این دو پارامتر
امکان مقایسه گیت‌ها را با محاسبه دستی فراهم می‌کنند
\cite[بخش \secnum{۴.۴}، صفحه ۱۵۵]{weste}
\cite[اسلایدهای ۳ تا ۵]{logicaleffortslides}.

\subsection{تأخیر نرمال‌شده و دو سهم آن}
\label{sec:delay-stage-effort}

\(\tau\) مقیاس زمانی مرجع برای یک فناوری و شرایط کاری مشخص است.
در تقریب \lr{RC} کتاب، اگر \(R\) مقاومت مؤثر \lr{Inverter}
واحد در قرارداد زمان نصف تغذیه و \(C\) ظرفیت ورودی یک ترانزیستور
واحد باشد، \(\tau\approx3RC\) است. ظرفیت ورودی خود
\lr{Inverter} مرجع \(3C\) بود. مقدار \(d\) با تقسیم تأخیر
زمانی \(t_{\mathrm{pd}}\) بر \(\tau\) بی‌بعد می‌شود. ظرفیت
ورودی مربوط به گذار را \(C_{\mathrm{in}}\) و ظرفیت خارجی‌ای
را که خروجی راه‌اندازی می‌کند \(C_L\) می‌نامیم.
کتاب و اسلایدها برای این بار خارجی از \(C_{\mathrm{out}}\)
استفاده می‌کنند. در این بخش \(C_L\) را به کار می‌بریم تا با
ظرفیت کل گره خروجی در مدل گذرا اشتباه نشود. روابط مدل چنین‌اند
\cite[روابط \secnum{۴.۲۰} تا \secnum{۴.۲۲}، صفحه ۱۵۵]{weste}
\cite[اسلاید ۴]{logicaleffortslides}.
\begin{equation}
  h=\frac{C_L}{C_{\mathrm{in}}},\qquad
  f=gh,\qquad
  d=\frac{t_{\mathrm{pd}}}{\tau}=f+p=gh+p.
  \label{eq:delay-normalized-gate}
\end{equation}
در این رابطه، \(h\) یا \lr{Electrical Effort} نسبت بار خارجی
به ظرفیت ورودی است. برای چند ورودی مشابه، \(h\) تنها زمانی برابر
تعداد آن‌هاست که ظرفیت هر یک با ظرفیت ورودی گیت راه‌انداز برابر
باشد و ظرفیت سیم قابل صرف‌نظر باشد. \(g\) یا \lr{Logical Effort}
اثر نوع گیت و ورودی موردنظر را بیان می‌کند. \(f\) یا
\lr{Stage Effort} سهم وابسته به بار است. \(p\) یا
\lr{Parasitic Delay} تأخیر نرمال‌شده در \(C_L=0\) است.
ظرفیت‌های درونی که در \(p\) منظور شده‌اند، در \(C_L\) دوباره
شمرده نمی‌شوند.

برای فهم \(g\)، گیت را با یک \lr{Inverter} مقایسه کنید که در
گذار موردنظر همان قدرت راه‌اندازی را دارد. اگر \(C_{\mathrm{in,ref}}\)
ظرفیت ورودی این \lr{Inverter} باشد، تعریف \(g\) چنین است
\cite[بخش \secnum{۴.۴.۱}، صفحه ۱۵۶]{weste}
\cite[اسلاید ۶]{logicaleffortslides}.
\begin{equation}
  g=\frac{C_{\mathrm{in}}}{C_{\mathrm{in,ref}}}.
  \label{eq:delay-gate-logical-effort}
\end{equation}
در تقریب \lr{RC}، مقاومت مؤثر گیت و \lr{Inverter} هم‌قدرت را
برابر و برابر \(R_{\mathrm{eq}}\) می‌گیریم. با تغییر یکنواخت اندازه
\lr{Inverter}، مقاومت و ظرفیت ورودی آن به‌طور معکوس تغییر
می‌کنند و \(R_{\mathrm{eq}}C_{\mathrm{in,ref}}\approx\tau\)
می‌ماند. بنابراین سهم تأخیر بار در واحد \(\tau\) از
\(R_{\mathrm{eq}}C_L/\tau\) به
\(C_L/C_{\mathrm{in,ref}}=gh\) می‌رسد. از این راه روشن می‌شود
که چرا \(g\) شیب نمودار تأخیر نسبت به \(h\) است.

\subsection{عرض ترانزیستورها و ضرایب مدل}
\label{sec:delay-logical-effort-sizing}

فرض کنیم مقاومت یک \lr{pMOS} واحد \(2R\) و مقاومت یک
\lr{nMOS} واحد \(R\) است. در این مدل \lr{Inverter} مرجع
با عرض‌های نسبی ۱ و ۲، در هر دو جهت مقاومت \(R\) و ظرفیت
ورودی \(3C\) دارد. برای \lr{NAND} دوورودی، در مسیر افت دو
\lr{nMOS} سری روشن‌اند. با عرض ۲ برای هرکدام، مقاومت مسیر
\(R/2+R/2=R\) می‌شود. در مسیر صعود، در حالت یک
\lr{pMOS} روشن، عرض ۲ کافی است. هر ورودی به یک
\lr{nMOS} و یک \lr{pMOS} با عرض ۲ وصل است، پس ظرفیت
هر ورودی \(4C\) و \(g=4/3\) است.

در \lr{NOR} دوورودی، دو \lr{pMOS} سری در مسیر صعود
قرار دارند. برای هرکدام عرض ۴ می‌گیریم تا
\(2R/4+2R/4=R\) شود. در مسیر افت، یک \lr{nMOS} با عرض ۱
برای مقاومت \(R\) کافی است. این بار هر ورودی ظرفیت \(5C\)
دارد و \(g=5/3\) می‌شود. با همین استدلال، برای \(n\) ورودی،
عرض ترانزیستورهای سری در \lr{NAND} برابر \(n\) و در
\lr{NOR} برابر \(2n\) است. جدول \ref{tab:delay-simple-gates}
ظرفیت هر ورودی و پارامترهای مدل را برای این گیت‌ها نشان می‌دهد
\cite[شکل \secnum{۴.۲۲} و جدول‌های \secnum{۴.۲}
و \secnum{۴.۳}، صفحات ۱۵۶ و ۱۵۷]{weste}
\cite[اسلایدهای ۶ تا ۹]{logicaleffortslides}.

\begin{table}[htbp]
  \centering
  \caption{پارامترهای گیت‌ها برای مقاومت راه‌اندازی یکسان.
  \(C\) ظرفیت ورودی ترانزیستور واحد و \(p_{\mathrm{inv}}\)
  تأخیر \lr{Parasitic} نرمال‌شده \lr{Inverter} است.}
  \label{tab:delay-simple-gates}
  \begin{tabular}{cccc}
    \toprule
    گیت & ظرفیت هر ورودی & \(g\) & \(p/p_{\mathrm{inv}}\) \\
    \midrule
    \lr{Inverter} & \(3C\) & \(1\) & \(1\) \\
    \lr{NAND} با \(n\) ورودی & \((n+2)C\) & \((n+2)/3\) & \(n\) \\
    \lr{NOR} با \(n\) ورودی & \((2n+1)C\) & \((2n+1)/3\) & \(n\) \\
    \bottomrule
  \end{tabular}
\end{table}

برای ستون آخر جدول، فقط ظرفیت‌های \lr{Diffusion} متصل به
خروجی را می‌شماریم. در \lr{NAND}، یک \lr{nMOS} با عرض
\(n\) و \(n\) ترانزیستور \lr{pMOS} با عرض ۲ به خروجی
متصل‌اند. مجموع عرض‌ها \(3n\) است. در \lr{NOR}، یک
\lr{pMOS} با عرض \(2n\) و \(n\) ترانزیستور \lr{nMOS}
با عرض ۱ باز هم مجموع \(3n\) را می‌دهند. با مقاومت مؤثر برابر
و با فرض ناحیه‌های \lr{Drain} مجزا، تأخیر بی‌بار تقریباً
\(n\) برابر مقدار \lr{Inverter} است. در تقریب \(3C\) برای
ظرفیت \lr{Diffusion} خروجی \lr{Inverter} واحد،
\(p_{\mathrm{inv}}\approx1\) است.

با بزرگ‌کردن یکنواخت یک گیت به اندازه \(k\)، مقاومت به
\(1/k\) مقدار قبلی می‌رسد و ظرفیت داخلی تقریباً \(k\) برابر
می‌شود. از این رو، حاصل‌ضرب آن‌ها و در نتیجه \(p\) در تقریب
اول تقریباً ثابت می‌ماند. گره‌های داخلی می‌توانند این تخمین را
تغییر دهند. برای نمونه، در گذار افت \lr{NAND} کتاب، ترانزیستور
پایینی آخر از همه روشن می‌شود و ظرفیت گره‌های داخلیِ از پیش
شارژشده نیز تخلیه می‌شود. با شماره‌گذاری از سمت زمین،
گره داخلی \(i\) ظرفیت \(nC\) دارد و با خروجی در مقاومت
\(iR/n\) مشترک است. با روش \lr{Elmore}، جمع
سهم خروجی و گره‌های داخلی چنین می‌شود
\cite[شکل \secnum{۴.۲۳} و رابطه \secnum{۴.۲۳}، صفحات ۱۵۷ و ۱۵۸]{weste}.
\[
  t_{\mathrm{pd}}\approx 3nRC+\sum_{i=1}^{n-1}\frac{iR}{n}(nC)
    =\left(3n+\frac{n(n-1)}{2}\right)RC.
\]
جمله دوم نشان می‌دهد چرا تخمین مبتنی بر ظرفیت خروجی به‌تنهایی،
برای پشته‌های سری بلند سهم مهمی از تأخیر را از دست می‌دهد.

\subsection{خواندن نمودار و محاسبه تأخیر \lr{FO4}}
\label{sec:delay-fo4}

در شکل \ref{fig:delay-gate-comparison}، عرض از مبدأ برابر \(p\)
و شیب برابر \(g\) است. با فرض \(p_{\mathrm{inv}}=1\)،
در \(h=3\) به‌ترتیب \(d=4\)، \(d=6\) و \(d=7\) برای
\lr{Inverter}، \lr{NAND} دوورودی و \lr{NOR} دوورودی
به دست می‌آید. مقایسه در \(h\) یکسان، یعنی مقایسه با نسبت بار
به ظرفیت ورودی یکسان. اگر ظرفیت ورودی گیت‌ها متفاوت باشد،
بارهای خارجیِ متناظر نیز متفاوت‌اند.

\begin{figure}[htbp]
  \centering
  \begin{latin}
  \begin{tikzpicture}[x=2.15cm,y=0.47cm]
    \draw[->] (0,0) -- (4.65,0) node[right] {\(h\)};
    \draw[->] (0,0) -- (0,9.05) node[above] {\(d\)};
    \node[below left] at (0,0) {\rl{۰}};
    \foreach \x/\ticklabel in {1/۱,2/۲,3/۳,4/۴} {
      \draw (\x,0.09) -- (\x,-0.09) node[below] {\rl{\ticklabel}};
    }
    \foreach \y/\ticklabel in {1/۱,2/۲,4/۴,6/۶,8/۸} {
      \draw (0.04,\y) -- (-0.04,\y) node[left] {\rl{\ticklabel}};
    }
    \draw[gray,densely dotted] (3,0) -- (3,7);
    \draw[thick] (0,1) -- (3.65,4.65) node[right] {\small Inverter};
    \draw[thick,blue!70!black,dashed] (0,2) -- (3.65,{2+3.65*4/3})
      node[right] {\small NAND2};
    \draw[thick,red!65!black,dashdotted] (0,2) -- (3.65,{2+3.65*5/3})
      node[right] {\small NOR2};
    \fill (3,4) circle (1.6pt);
    \fill[blue!70!black] (3,6) circle (1.6pt);
    \fill[red!65!black] (3,7) circle (1.6pt);
  \end{tikzpicture}
  \end{latin}
  \caption{تأخیر نرمال‌شده \(d\) در برابر \lr{Electrical Effort} بی‌بعد
  \(h\) برای سه گیت، با فرض \(p_{\mathrm{inv}}=1\).
  نمودار از رابطه \eqref{eq:delay-normalized-gate} و ضرایب
  جدول \ref{tab:delay-simple-gates} رسم شده است.}
  \label{fig:delay-gate-comparison}
\end{figure}

در \lr{Fanout-of-4 (FO4)}، یک \lr{Inverter} چهار ورودی
\lr{Inverter} هم‌اندازه را راه‌اندازی می‌کند. در مدل ساده،
\(g=1\)، \(h=4\) و \(p=1\) است. در نمونه کتاب برای یک فناوری
\qty{65}{\nano\metre} با
\(\tau=\qty{3}{\pico\second}\)، تأخیر زمانی چنین به دست می‌آید
\cite[مثال \secnum{۴.۱۰}، صفحه ۱۵۸]{weste}
\cite[اسلاید ۱۲]{logicaleffortslides}.
\[
  d_{\mathrm{FO4}}=1\cdot4+1=5,\qquad
  t_{\mathrm{FO4}}=d_{\mathrm{FO4}}\tau=\qty{15}{\pico\second}.
\]
مقدار ۵ نتیجه مدل نرمال‌شده است و مقدار زمانی به \(\tau\)
وابسته است. به همین دلیل، مقایسه برحسب \lr{FO4} باید با
مرجع همان فناوری و شرایط کاری انجام شود.

\subsection{دوره و فرکانس \lr{Ring Oscillator}}
\label{sec:delay-ring-oscillator}

در حلقه‌ای با تعداد فرد \(N\) از \lr{Inverter}ها، گذار پس
از یک دور با قطبیت مخالف و پس از دور دوم با قطبیت اولیه به
همان گره برمی‌گردد. هر گیت طی یک دوره، یک گذار صعودی و یک
گذار نزولی دارد. اگر گیت‌ها هم‌اندازه باشند و تأخیر هر گذار
به‌ترتیب \(t_\uparrow\) و \(t_\downarrow\) باشد، دوره \(T\) و
فرکانس \(f_{\mathrm{osc}}\) چنین‌اند
\cite[شکل \secnum{۴.۲۴} و مثال \secnum{۴.۱۱}، صفحات ۱۵۸ و ۱۵۹]{weste}
\cite[اسلایدهای ۱۰ و ۱۱]{logicaleffortslides}.
\begin{equation}
  T=N(t_\uparrow+t_\downarrow)=2N\bar t_{\mathrm{pd}},\qquad
  f_{\mathrm{osc}}=\frac{1}{T}.
  \label{eq:delay-ring-frequency}
\end{equation}
در این رابطه، \(\bar t_{\mathrm{pd}}\) میانگین تأخیر صعود و افت
یک مرحله است. اگر بار اضافی و ظرفیت سیم ناچیز باشند، هر
\lr{Inverter} یک ورودی هم‌اندازه را راه‌اندازی می‌کند. بنابراین
\(h=1\) و با فرض \(p=1\)، مقدار \(d=1\cdot1+1=2\) است.
با \(N=31\) و \(\tau=\qty{3}{\pico\second}\)، تأخیر هر
مرحله \(\bar t_{\mathrm{pd}}=\qty{6}{\pico\second}\)، دوره
\(T=\qty{372}{\pico\second}\) و فرکانس
\(f_{\mathrm{osc}}\approx\qty{2.7}{\giga\hertz}\) می‌شود.
عامل ۲ در دوره از دو دور لازم برای بازگشت به قطبیت اولیه می‌آید.
برای محاسبه فرکانس برحسب هرتز، دوره باید برحسب ثانیه باشد.

\subsection{قدرت راه‌اندازی، استخراج ضرایب و حدود مدل}
\label{sec:delay-characterization}

در کتابخانه \lr{Standard Cell}، برچسب‌هایی مثل \lr{X1} و
\lr{X4} اندازه راه‌اندازی را نسبت به یک مرجع کتابخانه بیان
می‌کنند. در مدل آموزشی، اگر \(C_{\mathrm{in,inv}}\) ظرفیت
ورودی \lr{Inverter} واحد باشد، قدرت راه‌اندازی نسبی \(x\)
با نسبت ظرفیت \lr{Inverter} هم‌قدرت به ظرفیت
\lr{Inverter} واحد تعریف می‌شود. با رابطه
\eqref{eq:delay-gate-logical-effort}، این تعریف و صورت دیگر
تأخیر چنین‌اند \cite[روابط \secnum{۴.۲۴} و
\secnum{۴.۲۵}، صفحه ۱۵۹]{weste}.
\begin{equation}
  x=\frac{C_{\mathrm{in}}}{gC_{\mathrm{in,inv}}},\qquad
  d=\frac{C_L/C_{\mathrm{in,inv}}}{x}+p.
  \label{eq:delay-gate-drive}
\end{equation}
برای نمونه، \lr{NAND} دوورودیِ طراحی‌شده با عرض‌های جدول،
ظرفیت هر ورودی \(4C\) و \(g=4/3\) دارد. در نتیجه با مرجع
\(C_{\mathrm{in,inv}}=3C\)، قدرت راه‌اندازی آن \(x=1\) است.
این برابری نشان می‌دهد که ظرفیت ورودی گیت‌های هم‌قدرت می‌تواند
متفاوت باشد.

برای استخراج پارامترها از داده اندازه‌گیری یا \lr{Datasheet}،
در شرایط یکسان خط \(t_{\mathrm{pd}}=a+b h\) را برازش می‌کنیم.
در این رابطه \(a\) عرض از مبدأ و \(b\) شیب، هر دو با واحد زمان‌اند.
اگر داده به شکل تأخیر برحسب ظرفیت بار ارائه شود، شیب زمانی
برحسب \(h\) با ضرب شیب تأخیر برحسب ظرفیت در \(C_{\mathrm{in}}\)
به دست می‌آید. برای \lr{Inverter}، \(g=1\) است و
\(\tau=b_{\mathrm{inv}}\). برای هر گیت، \(g=b/\tau\) و
\(p=a/\tau\) است. در مثال کتاب، شیب \lr{Inverter} واحد
\qty{12.4}{\pico\second} است. برای یکی از ورودی‌های
\lr{NAND} دوورودی، شیب \qty{15.5}{\pico\second} و عرض
از مبدأ \qty{25.4}{\pico\second} است و به ترتیب \(g\approx1.25\)
و \(p\approx2.05\) به دست می‌آید
\cite[بخش \secnum{۴.۴.۵} و روابط \secnum{۴.۲۶}
و \secnum{۴.۲۷}، صفحات ۱۵۹ و ۱۶۰]{weste}.

ضرایب مدل خطی برای بازه مشخصی از بار، شیب ورودی و شرایط کاری
مناسب‌اند. اگر ورودی کندتر شود، ترانزیستور دیرتر قدرت راه‌اندازی
کامل پیدا می‌کند. اگر چند ورودی هم‌زمان تغییر کنند، مسیرهای سری
و موازی می‌توانند تأخیر متفاوتی نسبت به گذار تک‌ورودی داشته باشند.
همچنین \lr{Velocity Saturation} و تغییر ولتاژ تغذیه و آستانه
بر تخمین تأخیر اثر می‌گذارند. در پشته سری، تغییر ولتاژ
\lr{Source} جریان لازم برای شارژ ظرفیت \(C_{gs}\) بین
\lr{Gate} و \lr{Source} را تغییر می‌دهد. ظرفیت \(C_{gd}\)
بین \lr{Gate} و \lr{Drain} نیز ورودی و خروجی را به هم
کوپل می‌کند. در نتیجه، تغییر ولتاژ یکی می‌تواند بار لازم برای
تغییر دیگری را افزایش دهد. هنگام استفاده از ضرایب
جدولی باید گذار فعال و فرض‌های اندازه‌گذاری آن‌ها را مشخص کرد
\cite[بخش \secnum{۴.۴.۶}، صفحات ۱۶۰ تا ۱۶۳]{weste}.

